\documentclass{article}
\usepackage[utf8]{inputenc}
\usepackage[T1]{fontenc}
\usepackage[french]{babel}
\usepackage{amsmath} % Pour les environnements mathématiques comme cases et la commande \pmod
\usepackage{enumitem} % Pour personnaliser les listes

\begin{document}

\textbf{Exercice 1 (3 points)}

Dans ce qui suit, $x$ et $y$ désignent des entiers.
Répondre par vrai ou faux en justifiant la réponse.

\begin{enumerate}[label=\alph*)]
    \item $x^3 \equiv x \pmod{2}$.
    \item Si $x \equiv 2 \pmod{14}$ alors $x \equiv 1 \pmod{7}$.
    \item Si $4x \equiv 10y \pmod{5}$ alors $x \equiv 0 \pmod{5}$.
    \item Si $\begin{cases} x \equiv 4 \pmod{5} \\ y \equiv 5 \pmod{8} \end{cases}$ alors $8x - 5y = 7$.
\end{enumerate}

\end{document}
